\documentclass[11pt]{article}
\usepackage{graphicx} % Required for inserting images
\setlength{\parindent}{0pt}
\usepackage[bookmarks=true, colorlinks=true, urlcolor=blue, linkcolor=black]{hyperref}

\usepackage{enumitem}
\usepackage[utf8]{inputenc}
\usepackage[T1]{fontenc}
\usepackage[english]{babel}
\usepackage{geometry}
\usepackage{lastpage}
\usepackage{fancyhdr}
\usepackage{datetime}
\usepackage{parskip}

% Link to other SOPs, forms
% Definitions for SOP sections 3.
\newcommand{\AdverseEvent}{\textbf{Adverse Event}: Any untoward or unfavorable medical occurrence in a human subject, including any abnormal sign, symptom, or disease, temporally associated with the subject's participation in the research, whether or not considered related to the subject's participation in the research.}

\newcommand{\ControlRoom}{\textbf{Control Room}: B17; The main room of the MRI suite that contains the computer displays and keyboards that control the MRI scanner. The Control Room is accessed from B15 that is controlled by an Electronic Lock.}

\newcommand{\DcmComputer}{\textbf{DICOM Computer}: The Mac Studio located in the Control Room the receives DICOM files from the scanner.}

\newcommand{\ElectronicLock}{\textbf{Electronic Lock}: All doors at the ND-HNC are controlled by electronic locks. The electronic lock system recognizes the owner's Irish1Card, logs the time of the attempted entry, and unlocks the door if the owner has been granted access.}

\newcommand{\EEG}{\textbf{EEG}: Electroencephalograpy, a technique used to measure changes in electromagnetic properties of neural tissue.}

\newcommand{\EquipmentRoom}{\textbf{Equipment Room}: B18; The room containing hardware for MRI operations.}

\newcommand{\FDA}{\textbf{FDA}: Food and Drug Administration.}

\newcommand{\GoodCatch}{\textbf{Good Catch}: Early recognition of an event with the potential to cause damage, injury, or illness.}

\newcommand{\Helper}{\textbf{Helper}: Second person in Zones 3 and 4 who assists in emergency responses, typically Level 2-3.}

\newcommand{\Incident}{\textbf{Incident}: An occurence, action, or event that warrants reporting (e.g. a policy is not followed and/or an individual's health is at risk). Includes Adverse Events, Near Misses, Unanticipated Problems, etc.}

\newcommand{\Lead}{\textbf{Lead}: Primary person in Zones 3 and 4 responsible for the Visitor's safety, in charge during emergency responses. Level 4 operator or MRI Technologist.}

\newcommand{\MagnetRoom}{\textbf{Magnet Room}: B17A; The room within the MRI suite that contains the MRI scanner. It is accessed from within the MRI Control room (B17) through a doorway that is controlled by an Electronic Lock.}

\newcommand{\NearMiss}{\textbf{Near Miss}: An incident that occurred but did not cause personal injury/illness e.g. a small/contained fire, trip on sidewalk or carpeting, falling object.}

\newcommand{\NDHNC}{\textbf{ND-HNC}: Notre Dame Human Neuroimaging Center, Lower Level of Veldman Family Psychology Clinic.}

\newcommand{\Operator}{\textbf{Operator}: MRI Technologist or Level 4 researcher, responsible for conducting research operations and maintaining equpiment. The Lead during emergency responses.}

\newcommand{\PIN}{\textbf{PIN}: Personal Identification Number, separate for Irish1Card and security.}

\newcommand{\ResearchAssistant}{\textbf{Research Assistant}: Individual with access Levels 1-3 to the MRI research suite. Typically the Helper during emergency responses.}

\newcommand{\SeriousInjury}{\textbf{Serious Injury}: An injury or illness that is life threatening; results in permanent impairment of a body function or permanent damage to a body structure; or necessitates medical or surgical intervention to preclude permanent damage or impairment.}

\newcommand{\StimControl}{\textbf{Stimulus Control Computer}: The Windows machines in the Control Room that deploy fMRI tasks and capture participant responses.}

% \newcommand{\Subject}{\textbf{Subject}: An individual who is participating in an experimental protocol at ND-HNC and has not taken the ND-HNC safety course.}

\newcommand{\Supervision}{\textbf{Supervision}: A supervisor is continuously present in both sight and sound range of the individual or task being monitored. This approach ensures the supervisor can actively observe, direct, and immediately intervene if a complication or safety issue arises.}

\newcommand{\TMS}{\textbf{TMS}: Transcranial magnetic stimulation, a technique used to disrupt normal operations of neural tissue.}

\newcommand{\TornadoWarn}{\textbf{Tornado Warning}: A tornado has been sighted or indicated by weather radar and there is imminent danger to life and property. Action is required.}

\newcommand{\TornadoWatch}{\textbf{Tornado Watch}: Tornados are possible in and near the watch area. Be ready to act quickly and move to the designated tornado safe area.}

\newcommand{\UnanticipatedProblem}{\textbf{Unanticipated Problem}: An incident, experience, or outcome that meets \textbf{all} of the following criteria: (1) unexpected (in terms of nature, severity, or frequency) given (a) the research procedures that are described in the protocol-related documents, such as the IRB-approved research protocol and informed consent document; and (b) the characteristics of the subject population being studied; (2) related or possibly related to participation in the research; and (3) suggests that the research places subjects or others at a greater risk of harm than was previously known or recognized.}

\newcommand{\VFPC}{\textbf{VFPC}: Veldman Family Psychology Clinic.}

\newcommand{\Visitor}{\textbf{Visitor}: An individual who has not taken the ND-HNC MRI Safety Course but has a legitimate reason for entering the MRI suites. Research participants and their legal guardians are Visitors. Other Visitors may include: individuals providing/receiving a tour of the VFPC and ND-HNC, faculty and staff who have not received MRI Safety Course, and research assistants.}



% figures and tables
\usepackage{graphicx}
\graphicspath{ {../../pics} }
\usepackage{rotating}
\usepackage{float}
\usepackage{subfig}
\usepackage{caption}
\DeclareCaptionLabelFormat{adja-page}{\hrulefill\\#1 #2 \emph{(previous page)}}


\renewcommand{\headrulewidth}{0pt}
% \newdateformat{yearmonth}{\THEYEAR-\THEMONTH-\THEDAY}
\pagestyle{fancy}
\fancyhead[L]{}
\fancyhead[R]{Approved: TBD}
\fancyfoot[C]{}
\fancyfoot[L]{ND-HNC SOP 201: v.0.1.0}
\fancyfoot[R]{p. \thepage /\pageref*{LastPage}}


% start document
\begin{document}

\renewcommand{\labelenumii}{\arabic{enumi}.\arabic{enumii}}
\renewcommand{\labelenumiii}{\arabic{enumi}.\arabic{enumii}.\arabic{enumiii}}
\renewcommand{\labelenumiv}{\arabic{enumi}.\arabic{enumii}.\arabic{enumiii}.\arabic{enumiv}}

\begin{center}
	\textbf{\Large SOP 201: MRI Safety Screen}\\
	\hrulefill
\end{center}


\section{Summary}
All individuals must be screened before they enter the Magnet Room. This document establishes MRI Safety screening procedures at the Notre Dame Human Neuroimaging Center (ND-HNC).


\section{Scope}
These policies apply to all individuals who enter Zones 3 \& 4 (B15, B17, B17a; Fig. \ref{fig:mri-zones}).


\section{Definitions}

\begin{itemize}
	\item \MagnetRoom
	\item \NDHNC
	\item \Operator
	\item \Visitor
\end{itemize}

\section{Notes}

\begin{itemize}
    \item Emergency response and maintenance personnel \textbf{CANNOT} enter the Magnet Room unless they have been screened and all metal objects have been removed.
	\item Emergency response personnel can \textbf{ONLY} enter Zone 3 without completing an MRI Safety screener if the door to the Magnet Room has been locked by key.
	\item ND-HNC only stores MRI Safety forms for staff and researchers. Visitor forms are maintained by their respective research teams.
\end{itemize}


\section{Policies \& Procedures}

The purpose of MRI Safety screening is to determine whether entry into the Magnet Room poses a health or safety risk. Accordingly, all individuals entering Zones 3 \& 4 must be screened according to the procedures below.

%
\subsection{MRI Prescreen}
\label{ssec:mri-pre}
Prescreening is conducted for Visitor's intending to enter the Magnet Room (e.g. a research participant). Before arriving at the ND-HNC, research personnel check for potential MRI contraindications with the prospective Vistor:

\begin{itemize}
    \item Non-MRI compatible medical implants.
    \begin{itemize}
        \item \textbf{Note}: Visitors with an implant must present documentation from their surgeon's office that the device is MRI compatible.
        \item \textbf{IUDs}: North American non-metallic intrauterine devices (e.g. Mirena) are approved. Metallic (e.g. Cu) and IUDs from other regions must be approved by the MRI Technologist.
    \end{itemize}
    \item Other metal embedded in the body (e.g. shrapnel, unremovable piercings).
    \item Pregnancy.
    \item Claustrophobia (the Mock can be used to assess comfort levels).
    \item Tattoo on face or neck, including permanent makeup.
\end{itemize}

%
\subsection{MRI Preparation}
\label{ssec:mri-prep}
Visitors intending to enter the Magnet Room (e.g. a research participant) should be informed of the following restrictions, before arriving at the ND-HNC, to guide their preparations:

\begin{itemize}
    \item Clothing should be free of metallic objects and clasps; bras should not have an underwire. Anti-odorant materials (e.g. Lululemon's Silverescent) are not MRI compatible.
    \item Jewelry will need to be removed, including piercings if they are metallic.
    \item Cosmetics such as colored contacts, magnetic lashes, and makeup or nail polish with glitter will have to be removed. Tanning lotions are not MRI compatible.
    \item Hair must be free of all accessories (metallic clips, barrettes, scrunchees, etc.).
    \item If Visitors have corrected vision, they should wear contacts or know their prescription.
\end{itemize}

%
\subsection{MRI Safety Screen}
\label{ssec:mri-safety}
All who enter Zone 3, regardless of whether they will enter Zone 4 (Magnet Room) fill out Form 001: MRI Safety. ND-HNC staff and research personnel submit a form (a) annually, and (b) whenever a change occurs in their MRI compatibility. Visitors fill out the form in Zone 1, and researchers with at least Level 2 access should be available to answer any questions.

\begin{itemize}
    \item For Visitors, Form 001 must be filled out on the same day as entering Zone 3.
    \item Each question/prompt must be answered individually.
    \item In the case of minors, their legal guardian must fill out the form.
    \begin{itemize}
        \item Note: verify guardianship, which may differ from the individual who brought the minor to the center.
    \end{itemize}
    \item The Visitor signs and dates the form in the designated location.
\end{itemize}


\subsection{MRI Safety Screen, Review 1}
\label{ssec:mri-rev1}
A researcher with at least Level 2 access will review the Visitor's responses on Form 001.

\begin{itemize}
    \item Discuss any potential contraindications, points of confusion, or concerns. Inform the Operator about any items discussed or addressed.
    \item After reviewing Form 001, the researcher prints and signs their name in the designated location.
    \item If the researcher determines the Visitor is MRI compatible, they may escort the Visitor into Zones 2 \& 3.
\end{itemize}

% Screen Two
\subsection{MRI Safety Screen, Review 2}
\label{ssec:mri-rev2}
A second review is required for Visitors entering Zone 4 (Magnet Room).

\begin{itemize}
    \item The Operator conducts the second review, who must not be the same person that conducted Review 1.
    \item Any potential contraindications, points of confusion, or concerns are addressed, including those raise by the first reviewer
    \item The Operator signs and dates Form 001.
    \item The Visitor may enter Zone 4 only after the Operator has determined that it is safe to do so.
\end{itemize}


\section{Figures}

% MRI zones
\begin{figure}[H]
	\centering
	\includegraphics[width=\textwidth]{mri_zones.jpg}
	\caption{Schematic of ND-HNC MRI Zones. \textbf{Orange} Zone 1. \textbf{Purple} Zone 2. \textbf{Green} Zone 3. \textbf{Blue} Zone 4.}
	\label{fig:mri-zones}
\end{figure}


\end{document}